\section{EXAMPLE ~\ref{example:counter} EXPLICIT CALCULATION }

Calculation for the~\Cref{example:counter} in detail.\\
We take a three-state, two action Markov decision process with initial state $\emptyset$ and a uniform prior policy:
\[
\pi(a_1|\emptyset) = \pi(a_2|\emptyset) = \frac{1}{2}
\]
with transition dynamics depicted on the diagram below:
% https://q.uiver.app/#q=WzAsNSxbMiwwLCJcXGVtcHR5c2V0Il0sWzMsMCwic18yIl0sWzEsMCwic18xIl0sWzQsMCwiXFxtYXRoYmJ7UH0oUnxzXzIpID0gMSJdLFswLDAsIlxcbWF0aGJie1B9KFJ8c18xKSA9IDEiXSxbMCwxLCJcXGZyYWN7MX17NH0iLDAseyJjdXJ2ZSI6LTIsImNvbG91ciI6WzAsNjAsNjBdLCJzdHlsZSI6eyJib2R5Ijp7Im5hbWUiOiJkYXNoZWQifX19LFswLDYwLDYwLDFdXSxbMCwyLCJcXGZyYWN7M317NH0iLDIseyJjdXJ2ZSI6MiwiY29sb3VyIjpbMCw2MCw2MF0sInN0eWxlIjp7ImJvZHkiOnsibmFtZSI6ImRhc2hlZCJ9fX0sWzAsNjAsNjAsMV1dLFswLDEsIlxcZnJhY3sxfXsyfSIsMix7ImN1cnZlIjoyfV0sWzAsMiwiXFxmcmFjezF9ezJ9IiwwLHsiY3VydmUiOi0yfV1d
\[\begin{tikzcd}
	{s_1} & \emptyset & {s_2}
	\arrow["{\frac{3}{4}}"', color={rgb,255:red,214;green,92;blue,92}, curve={height=12pt}, dashed, from=1-2, to=1-1]
	\arrow["{\frac{1}{2}}", curve={height=-12pt}, from=1-2, to=1-1]
	\arrow["{\frac{1}{4}}", color={rgb,255:red,214;green,92;blue,92}, curve={height=-12pt}, dashed, from=1-2, to=1-3]
	\arrow["{\frac{1}{2}}"', curve={height=12pt}, from=1-2, to=1-3]
\end{tikzcd}\]
where red dotted lines correspond to action $a_1$, and solid black lines to action $a_2$. We also assume that:
\[
r(s_0, \cdot) = \log 1 \qquad r(s_1, \cdot) = \log \frac{1}{3} \qquad r(s_2, \cdot) = \log \frac{2}{3}
\]
\subsection{Control-as-inference operator}
The first iteration of $\pi^\mathcal{F}_1(\cdot|\emptyset)$ computes the probability of getting reward for action $a_1$:
\[
\pi^\mathcal{F}_1(a_1|\emptyset) = \mathbb{P}(a_1|\mathcal{O}^2_{1:2}, \emptyset) = \frac{\mathbb{P}(\mathcal{O}^2_{1:2}|a_1, \emptyset)\mathbb{P}(a_1|\emptyset)}{\mathbb{P}(\mathcal{O}^2_{1:2})}
= \left(\frac{1}{2}\cdot \frac{3}{4}\cdot\frac{1}{3} + \frac{1}{2}\cdot \frac{1}{4}\cdot\frac{2}{3}\right)/{\mathbb{P}(\mathcal{O}^2_{1:2})}
= \frac{5}{24} / {\mathbb{P}(\mathcal{O}^2_{1:2})}
\]
and for $a_2$:
\[
\pi^\mathcal{F}_1(a_2|\emptyset) = \mathbb{P}(a_2|\mathcal{O}^2_{1:2}, \emptyset) = \frac{\mathbb{P}(\mathcal{O}^2_{1:2}|a_2, \emptyset)\mathbb{P}(a_2|\emptyset)}{\mathbb{P}(\mathcal{O}^2_{1:2})}
= \left(\frac{1}{2}\cdot \frac{1}{2}\cdot\frac{1}{3} + \frac{1}{2}\cdot \frac{1}{2}\cdot\frac{2}{3}\right)/{\mathbb{P}(\mathcal{O}^2_{1:2})}
= \frac{6}{24} / {\mathbb{P}(\mathcal{O}^2_{1:2})}
\]
After normalisation, we get:
\[
\pi^\mathcal{F}_1(a_1|\emptyset) = \frac{5}{11} \qquad \pi^\mathcal{F}_1(a_2|\emptyset) = \frac{6}{11}
\]
Now, applying the same computation second time - for $a_1$:
\[
\pi^\mathcal{F}_2(a_1|\emptyset) = \mathbb{P}(a_1|\mathcal{O}^2_{1:2}, \emptyset) = \frac{\mathbb{P}(\mathcal{O}^2_{1:2}|a_1, \emptyset)\pi^\mathcal{F}_1(a_1|\emptyset)}{\mathbb{P}(\mathcal{O}^2_{1:2})}
= \left(\frac{5}{11}\cdot \frac{3}{4}\cdot\frac{1}{3} + \frac{5}{11}\cdot \frac{1}{4}\cdot\frac{2}{3}\right)/{\mathbb{P}(\mathcal{O}^2_{1:2})}
= \frac{25}{11\cdot 12} / {\mathbb{P}(\mathcal{O}^2_{1:2})}
\]
and for $a_2$:
\[
\pi^\mathcal{F}_2(a_2|\emptyset) = \mathbb{P}(a_1|\mathcal{O}^2_{1:2}, \emptyset) = \frac{\mathbb{P}(\mathcal{O}^2_{1:2}|a_2, \emptyset)\pi^\mathcal{F}_1(a_2|\emptyset)}{\mathbb{P}(\mathcal{O}^2_{1:2})}
= \left(\frac{6}{11}\cdot \frac{3}{4}\cdot\frac{1}{3} + \frac{6}{11}\cdot \frac{1}{4}\cdot\frac{2}{3}\right)/{\mathbb{P}(\mathcal{O}^2_{1:2})}
= \frac{36}{11\cdot 12} / {\mathbb{P}(\mathcal{O}^2_{1:2})}
\]
Again applying normalisation:
\[
\pi^\mathcal{F}_1(a_1|\emptyset) = \frac{25}{61} \qquad \pi^\mathcal{F}_1(a_2|\emptyset) = \frac{36}{61}
\]
\subsection{Raising temperature}
Raising temperature policy $\pi_{\alpha(2)}(\cdot|\emptyset)$ first modifies the MDP by setting the rewards:
\[
r_2(s_1) = \log \frac{1}{9} \qquad r_2(s_2) = \log \frac{4}{9}
\]
and then recomputes the posterior for $a_1$:
\[
\pi_{\alpha(2)}(a_1|\emptyset) = \mathbb{P}(a_1|\mathcal{O}^2_{1:2}, \emptyset) = \frac{\mathbb{P}(\mathcal{O}^2_{1:2}|a_1, \emptyset)\mathbb{P}(a_1|\emptyset)}{\mathbb{P}(\mathcal{O}^2_{1:2})}
= \left(\frac{1}{2}\cdot \frac{3}{4}\cdot\frac{1}{9} + \frac{1}{2}\cdot \frac{1}{4}\cdot\frac{4}{9}\right)/{\mathbb{P}(\mathcal{O}^2_{1:2})}
= \frac{7}{72} / {\mathbb{P}(\mathcal{O}^2_{1:2})}
\]
and $a_2$:
\[
\pi_{\alpha(2)}(a_2|\emptyset) = \mathbb{P}(a_1|\mathcal{O}^2_{1:2}, \emptyset) = \frac{\mathbb{P}(\mathcal{O}^2_{1:2}|a_2, \emptyset)\mathbb{P}(a_2|\emptyset)}{\mathbb{P}(\mathcal{O}^2_{1:2})}
= \left(\frac{1}{2}\cdot \frac{1}{2}\cdot\frac{1}{9} + \frac{1}{2}\cdot \frac{1}{2}\cdot\frac{4}{9}\right)/{\mathbb{P}(\mathcal{O}^2_{1:2})} = \frac{10}{72}/{\mathbb{P}(\mathcal{O}^2_{1:2})}
\]
Again applying normalisation:
\[
\pi_{\alpha(2)}(a_1|\emptyset) = \frac{7}{17} \qquad \pi_{\alpha(2)}(a_2|\emptyset) = \frac{10}{17}
\]

